\begin{enumerate}[label=\thesubsection.\arabic*, ref=\thesubsection.\theenumi]
\item For a non-singular $n \times n$ matrix $\vec{A}$,
\begin{align}
	\vec{A}\,\text{adj}\,\vec{A} &= \mydet{\vec{A}}\vec{I}, \quad
	\vec{A}^{-1} = \frac{\text{adj}\,\vec{A}}{\mydet{\vec{A}}}
	\label{eq:EE26BTECH11202-adj}
	\\
	\mydet{\text{adj}\,\vec{A}} &= \mydet{\vec{A}}^{n-1}
	\label{eq:EE26BTECH11202-adj-det}
\end{align}
\solution
Taking the determinant on both sides of the first equation in
\eqref{eq:EE26BTECH11202-adj},
\begin{align}
	\mydet{\vec{A}}\mydet{\text{adj}\,\vec{A}} = \mydet{\mydet{\vec{A}}\vec{I}} = \mydet{\vec{A}}^{n}
\end{align}
which yields \eqref{eq:EE26BTECH11202-adj-det} when $\mydet{\vec{A}} \ne 0$.
\item The linear system $\vec{A}\vec{x} = \vec{b}$ with $\mydet{\vec{A}} \ne 0$ has the unique solution
\begin{align}
	\vec{x} = \vec{A}^{-1}\vec{b}
	\label{eq:EE26BTECH11202-linsys}
\end{align}
\item The area of the triangle with vertices $\vec{A}, \vec{B}, \vec{C}$ is
\begin{align}
	\frac{1}{2}\norm{\brak{\vec{B}-\vec{A}}\times\brak{\vec{C}-\vec{A}}}
	\label{eq:EE26BTECH11202-area}
\end{align}
In $\mathbb{R}^2$ this reduces to $\frac{1}{2}\abs{\mydet{\vec{B}-\vec{A} & \vec{C}-\vec{A}}}$ and the points are collinear if and only if
\begin{align}
	\text{rank}\myvec{\vec{B}-\vec{A} & \vec{C}-\vec{A}} = 1
	\iff
	\mydet{\vec{B}-\vec{A} & \vec{C}-\vec{A}} = 0
	\label{eq:EE26BTECH11202-collinear}
\end{align}
\item For any vectors $\vec{a}, \vec{b}$,
\begin{align}
	\norm{\vec{a}-\vec{b}}^2 &= \norm{\vec{a}}^2 + \norm{\vec{b}}^2 - 2\vec{a}^{\top}\vec{b}
	\label{eq:EE26BTECH11202-norm-id}
	\\
	\vec{a}^{\top}\vec{b} &= 0 \iff \vec{a} \perp \vec{b}
	\label{eq:EE26BTECH11202-ortho}
\end{align}
\item The angle $\theta$ between the lines $\vec{x} = \vec{A}_1 + \lambda\vec{m}_1$ and $\vec{x} = \vec{A}_2 + \mu\vec{m}_2$ is given by
\begin{align}
	\cos\theta = \frac{\vec{m}_1^{\top}\vec{m}_2}{\norm{\vec{m}_1}\norm{\vec{m}_2}}
	\label{eq:EE26BTECH11202-angle}
\end{align}
\item The foot of the perpendicular from $\vec{Q}$ to the line $\vec{x} = \vec{P}_0 + \lambda\vec{m}$ is $\vec{F} = \vec{P}_0 + \lambda_0\vec{m}$, where
\begin{align}
	\lambda_0 = \frac{\vec{m}^{\top}\brak{\vec{Q}-\vec{P}_0}}{\vec{m}^{\top}\vec{m}}
	\label{eq:EE26BTECH11202-foot}
\end{align}
\solution
$\vec{F}-\vec{Q}$ is orthogonal to $\vec{m}$. From \eqref{eq:EE26BTECH11202-ortho},
\begin{align}
	\vec{m}^{\top}\brak{\vec{P}_0 + \lambda_0\vec{m} - \vec{Q}} = 0
	\implies \lambda_0 = \frac{\vec{m}^{\top}\brak{\vec{Q}-\vec{P}_0}}{\vec{m}^{\top}\vec{m}}
\end{align}
\item The shortest distance between the skew lines $\vec{x} = \vec{A}_1 + \lambda\vec{m}_1$ and $\vec{x} = \vec{A}_2 + \mu\vec{m}_2$ is
\begin{align}
	d = \frac{\abs{\brak{\vec{A}_2-\vec{A}_1}^{\top}\brak{\vec{m}_1\times\vec{m}_2}}}{\norm{\vec{m}_1\times\vec{m}_2}}
	\label{eq:EE26BTECH11202-skew}
\end{align}
\item For events $A, B$ and a partition $A_1, A_2, A_3$ of the sample space with $\pr{A_i}=p_i$, $\pr{E\vert A_i}=q_i$,
\begin{align}
	\pr{A\vert B} &= \frac{\pr{AB}}{\pr{B}}, \quad \pr{A+B} = \pr{A}+\pr{B}-\pr{AB}
	\label{eq:EE26BTECH11202-cond}
	\\
	\pr{E} &= \sum_{i=1}^{3} p_i q_i = \vec{p}^{\top}\vec{q}
	\label{eq:EE26BTECH11202-total}
	\\
	\pr{A_j\vert E} &= \frac{p_j q_j}{\vec{p}^{\top}\vec{q}}
	\label{eq:EE26BTECH11202-bayes}
\end{align}
and $A, B$ are independent if and only if $\pr{AB} = \pr{A}\pr{B}$.
\item The conic $\vec{x}^{\top}\vec{V}\vec{x} = c$, with $\vec{V}$ symmetric positive definite having eigenvalues $\lambda_1, \lambda_2$, is an ellipse with semi-axes $\sqrt{c/\lambda_1}$ and $\sqrt{c/\lambda_2}$. Its area is
\begin{align}
	\frac{\pi c}{\sqrt{\mydet{\vec{V}}}}
	\label{eq:EE26BTECH11202-ellipse}
\end{align}
\item A linear programming problem is
\begin{align}
	\max_{\vec{x}} \; \vec{c}^{\top}\vec{x} \quad \text{s.t.} \quad
	\vec{A}\vec{x} \preceq \vec{b}, \; \vec{x} \succeq \vec{0}
	\label{eq:EE26BTECH11202-lpp}
\end{align}
and, if the feasible region is a bounded polygon, the optimum is attained at a corner point. If two corner points give the same optimum, so does every point on the segment joining them.
\item Standard integrals,
\begin{align}
	\int \frac{dx}{\sqrt{x^2+a^2}} &= \log\abs{x+\sqrt{x^2+a^2}} + C
	\label{eq:EE26BTECH11202-int-root}
	\\
	\int \sec x\, dx &= \log\abs{\sec x + \tan x} + C
	\label{eq:EE26BTECH11202-int-sec}
\end{align}
\item A differentiable $f$ is increasing on an interval if $f^{\prime}\brak{x} > 0$ there. For $y=u^{v}$,
\begin{align}
	\diff{y}{x} = y\,\diff{}{x}\brak{v\log u}
	\label{eq:EE26BTECH11202-logdiff}
\end{align}
\item For finite sets $A, B$ with $\abs{A}=m$, $\abs{B}=n$,
the number of relations from $A$ to $B$ is $2^{mn}$. A bijection $A \to B$ exists only if $m = n$.
\end{enumerate}
